\chapter{Overview}
\label{chap:overview}

The goal is the \v{S}varc--Milnor lemma in a form that is short to prove, yields the standard
corollaries (finite generation, change of generators, finite-index subgroups,
$\Z^n \sim_{QI} \R^n$) directly, and provides the quasi-isometry layer that later work on
hyperbolic groups needs. The reference text is C.~L\"oh, \emph{Geometric Group Theory: An
Introduction}, Springer, 2017; item numbers below refer to it.

\section*{Statement}

Let a group $G$ act by isometries on a pseudometric space $X$. Assume that $X$ is
\emph{coarsely geodesic}: there are constants $r, L, C$ such that any two points $x, y$ are joined
by a chain $x = x_0, x_1, \dots, x_n = y$ with $d(x_i, x_{i+1}) \le r$ and $n \le L\,d(x, y) + C$.
Assume that the action is \emph{cobounded}: some bounded $s \subseteq X$ has translates covering
$X$. Fix $x_0 \in s$ and put $S = \{g \in G \mid \exists a, b \in s,\ d(a, g b) \le r\}$. Then
\begin{enumerate}
  \item $S$ generates $G$;
  \item $d(g x_0, h x_0) \le (r + 2 \diam s)\,|g^{-1}h|_S$ and
        $|g^{-1}h|_S \le L\,d(g x_0, h x_0) + C$ for all $g, h \in G$;
  \item every point of $X$ is within $\diam s$ of the orbit $G x_0$;
  \item if moreover the action is \emph{metrically proper}, then $S$ is finite, so $G$ is
        finitely generated, and the orbit map is a quasi-isometry for the word metric of any
        finite generating family.
\end{enumerate}
Parts 1--3 need no properness. Part 1 uses only the existence of $r$-chains; the bound on their
length is needed for the lower bound in part 2 and cannot be dropped: $\Z$ acting by
translations on $(\R, \sqrt{d})$ satisfies every other hypothesis, but its orbit map is not a
quasi-isometry.

\section*{Why chains}

The chain hypothesis replaces quasi-geodesic paths. It makes the proof a finite counting
argument, needs no theory of geodesic spaces, and holds for word metrics themselves ($r = L = 1$,
$C = 0$), so the corollaries follow by letting a group act on itself. No generality is lost: in
every pseudometric space the chain condition is equivalent, with other constants, to L\"oh's
(non-continuous) quasi-geodesic condition (Definition 5.3.4).

\section*{Layout}

Chapter~\ref{chap:quasi-isometry} develops quasi-isometries,
Chapter~\ref{chap:word-metric} word metrics, Chapter~\ref{chap:coarsely-geodesic} coarsely
geodesic spaces and Chapter~\ref{chap:actions} the hypotheses on the action.
Chapter~\ref{chap:svarc-milnor} proves the lemma and Chapter~\ref{chap:corollaries} its
corollaries. Chapter~\ref{chap:coarse-spaces} records the coarse-space library the project
contains and its relation to the metric notions.
